\documentclass[../../../include/open-logic-section]{subfiles}

\begin{document}

\olfileid{cmp}{thy}{nou}
\olsection{不存在通用可计算函数}

% Part: computability
% Chapter: computability-theory
% Section: no-universal-function

虽然存在一个部分可计算函数，它对所有部分可计算函数而言是全函数，但不存在一个全可计算函数能作为所有全可计算函数的通用函数。

\begin{thm}
\ollabel{thm:no-univ}
不存在通用可计算函数。换言之，任何满足以下条件的函数 $\fn{Un}'(k, x)$ 都不是可计算的：如果 $f(x)$ 是全可计算函数，则存在自然数~$k$ 使得对每个~$x$ 都有 $f(x) = \fn{Un}'(k,x)$。
\end{thm}

\begin{proof}
证明采用简单的对角线法。如果 $\fn{Un}'(k,x)$ 是全的且可计算的，那么
\[
d(x) = \fn{Un}'(x, x) + 1
\]
也将是全的且可计算的。然而，根据定义，$d(k)$ 不等于 $\fn{Un}'(k,k)$。因此，对每个 $k$，$d(x)$ 和~$\fn{Un}'(k, x)$ 至少在一个点（即 $x = k$）处取值不同。
\end{proof}

\begin{explain}
上文的 \olref[uni]{thm:univ-comp} 表明，我们可以避开这个对角线论证，但代价是允许通用函数为部分函数。也就是说，$\fn{Un}$ 对于全可计算函数是通用的，只是它本身不是全函数。对角线论证在部分函数的情形下失效。
\end{explain}

\begin{prob}
为了理解为什么 \olref{thm:no-univ} 证明中的对角线论证在部分函数情形下失效，考虑函数 $f(x) \simeq \fn{Un}(x,x)+1$。它是部分可计算的吗？如果是，它就有一个指标~$e$，即 $f(x) \simeq
  \fn{Un}(e,x)$。关于~$f(e)$ 你能得出什么结论？
\end{prob}

\end{document}